% Short running form: the full title is wider than the box the running head
% leaves beside the Sunday's name.
\runninghead{The Publicans Go Before You}
\properlane{two-peoples}{The Publicans Go Before You: Two Sons and Two Peoples}
\label{sec:the-publicans-go-before-you}

Jerome, John Chrysostom and Thomas Aquinas read the two sons as two peoples.
In their reading the first son is the people of the nations, who refused God's
call and afterwards, at the Saviour's coming, repented and worked in the
vineyard; the second is the people of Israel, who answered at Sinai that they
would do all that the Lord had spoken, and did not go. The parable's own
application is narrower and sharper. It sets the publicans and harlots who
believed John before the chief priests and ancients of the people who did not.
Jerome finds the same two peoples in Ezekiel's two cases, and he hears the
prophet's call to a new heart addressed \latin{ad utrumque populum}, to both.
Chrysostom hears in ``go before you'' a hope held out to those who are behind.
Hilary of Poitiers assigns the sons the other way round, and his reading stands
beside theirs.

\subsection{A judgment drawn from their own mouths}

The parable is spoken to men who had just declined to answer a question. Asked
whether John's baptism was from heaven or from men, the chief priests and
ancients said, ``We know not'' (Mt 21:27). Jesus then tells a story and lets
them pass sentence on it. Chrysostom watches the method: ``Again He convicts
them by a parable, intimating both their unreasonable obstinacy, and the
submissiveness of those who were utterly condemned by them.'' He causes ``the
judgment to be delivered by themselves'', so that they are ``self-condemned''.
Had he said at once that harlots go before them, ``the word would have seemed
to them to be offensive; but now, being uttered after their own judgment it
appears to be not too hard.'' The charge that follows is exact. John came
``unto you'', not to the publicans, and came ``in the way of righteousness'', so
that no fault could be found with the messenger; the publicans heeded him; and
the chief priests did not heed him even after the publicans had done so. ``To
you he came, and ye accepted him not; he came not to them, and they receive
him, and not even them did ye take for instructors'' (\work{Homilies on
Matthew} 67.2--3).

The saying that closes the parable is not, for Chrysostom, a closed door. ``The
word, `go before you,' is not as though these were following, but as having a
hope, if they were willing. For nothing, so much as jealousy, rouses the
grosser sort.'' The publicans and harlots are named so that the chief priests
might be provoked to follow them (67.3). The parable's hearers are therefore
the religious leaders of Jerusalem who would not believe John, and the
parable's purpose, in Chrysostom's reading, is still their conversion.

\subsection{The nations and the people of Sinai}

The three witnesses of this interpretation widen the frame from the hearers to
the peoples. Chrysostom states it directly: ``these two children declare what
came to pass with respect to both the Gentiles and the Jews. For the former not
having undertaken to obey, neither having become hearers of the law, showed
forth their obedience in their works; and the latter having said, `All that
the Lord shall speak, we will do, and will hearken,' in their works were
disobedient.'' The words of the second son are the words of the people at
Sinai (Exod 19:8; 24:3, 7), and Chrysostom adds Paul's sentence so that ``they
might not think the law would benefit them'': ``Not the hearers of the law are
just before God, but the doers of the law shall be justified'' (Rom 2:13;
\work{Homilies on Matthew} 67.2).

Jerome's commentary gives the same identification with its reasons. He first
calls the two sons those of Luke's parable, the frugal son and the prodigal,
\latin{frugi et luxuriosus}. Then: \latin{Primo dicitur Gentilium populo per
naturalis legis notitiam: Vade, et operare in vinea mea}, to the first, the
people of the Gentiles, it is said through its knowledge of the natural law,
Go and work in my vineyard; and he gives the content of that law from Tobias,
not to do to another what one would not have done to oneself. \latin{Qui
superbe respondit: Nolo. Postea vero in adventu Salvatoris, acta poenitentia,
operatus est in vinea Dei, et sermonis contumaciam labore correxit}. That
people answered proudly, I will not; but afterwards, at the Saviour's coming,
it did penance, worked in God's vineyard, and corrected the defiance of its
word by its labour. \latin{Secundus autem filius, populus Iudaeorum est, qui
respondit Moysi: Omnia quaecumque dixerit Dominus faciemus \ldots\ et non ivit
in vineam, quia, interfecto patrisfamilias filio, se putavit haeredem}. The
second son is the people of the Jews, who answered Moses, All that the Lord
shall say we will do, and did not go into the vineyard, because, having killed
the householder's son, it thought itself the heir (\work{Commentarii in
Matthaeum} III, PL 26, 155--156). Jerome's last clause borrows the words of the
parable that follows, where the husbandmen say of the son, ``This is the heir:
come, let us kill him'' (Mt 21:38). On v.~31 he says that the publicans and
harlots, who had refused God by their evil works, afterwards received John's
baptism of penance.

Aquinas gives the reading first among his three: \latin{Iste homo, Deus est;
duo filii sunt duo populi}, this man is God, the two sons are two peoples; and
\latin{Primus est populus gentilium, qui incepit a Noe, sicut populus Iudaeorum
ab Abraham}, the first is the people of the Gentiles, which began from Noah as
the people of the Jews began from Abraham. The second son professes justice, as
the people said at Sinai, \latin{Omnia quaecumque praeceperit Dominus
faciemus}. Then he does what neither Jerome nor Chrysostom does here: he turns
the second son on the Church. Clerics and religious, too, profess justice, and
a promise made and not kept is a double sin, of disobedience and of a broken
vow (\work{Super Matthaeum} c.~21, pp.~275--276).

The three agree on the identification and on its ground, the promise given at
Sinai, set against nations that had made no such promise. They differ in what they draw from
it. Jerome ties the second son to the killing of the heir; Chrysostom keeps
open a hope for those who were behind; Aquinas applies the second son's fault to
the professed members of the Church.

\subsection{The way the house of Israel calls unequal}

Jerome reads the first reading in the same key. After stating the principle
that each is judged where he is found, he offers a second understanding of
Ezekiel's two cases: \latin{Potest et hoc intelligi: Iustus prius populus
Israel, avertit se a iustitia sua, quia iustitiae reliquit auctorem, et fecit
iniquitatem, Dei Filium denegando}; it can be understood thus, that the people
of Israel, once just, turned from its justice, because it left the author of
justice and did iniquity in denying the Son of God. Conversely, \latin{gentium
populus non habens notitiam Dei, et impius, si se averterit ab impietate sua,
quam prius operatus est in idololatria \ldots\ ipse prius mortuus vivificabit
animam suam}: the people of the nations, not knowing God and impious, if it
turns from the impiety it wrought in idolatry, though once dead, will give its
soul life. On the complaint of v.~29 Jerome writes a hard sentence of his own:
\latin{Usque hodie Israel blasphemat Deum, cur populum suum reliquerit, et
gentium assumpserit multitudinem}, to this day Israel blasphemes God for
leaving his people and taking up the multitude of the nations; and he calls the
sentence just, citing Luke's form of the husbandmen, whose vineyard is given to
others (Lk 20:16). The same commentary does not end there. On v.~30 he says
that whether one is of the nations or of Israel, \latin{Non est personarum
acceptio apud Deum \ldots\ sed unusquisque sua coronabitur fide}, there is no
respect of persons with God, but each will be crowned by his own faith. And on
the call to a new heart in v.~31, though it is said properly to Israel,
\latin{Potest autem et ad utrumque populum intelligi: ut et Israel, et gentium
turba, sua vitia derelinquens, convertatur ad eum qui possit sanare
contritiones suas}, it can be understood of both peoples, that Israel and the
crowd of the nations alike, leaving their vices, may turn to the one who can
heal their wounds. \latin{Cor novum Israelis, est credere in eum quem prius
negaverat. Cor novum gentium, est idola derelinquere}: Israel's new heart is to
believe in him whom it had denied; the nations' new heart is to leave their
idols (\work{In Hiezechielem} VI, PL 25, 181).

Augustine hears Ezekiel's complaint in another place and names no people. In
the psalm verse ``with the froward thou wilt be froward'' he finds the
complaint of those who say, ``The way of the Lord is not right'': ``Now this
seems froward to the froward, that Thou wilt make them whole that confess
their sins.'' God humbles those who are ``ignorant of God's righteousness, and
seek to establish their own'' (\work{Enarrationes in Psalmos} 17.27--28). For
Augustine the complaint belongs to anyone who resents the mercy shown to the
penitent; Jerome gives it to a people.

\subsection{Hilary: the sons the other way round}

Hilary of Poitiers read the parable of the same two groups and assigned them in
the opposite order. His text, with the Lectionary's order of the sons, had the
Pharisees answer by naming the younger, and he begins from the difficulty that
the ordinary identification meets. The elder son refused and repented;
\latin{Atquin Israel non poenituit, sed in Dominum manus intulit}, yet Israel
did not repent but laid hands on the Lord. The younger promised and did not go;
yet \latin{gentium peccatorumque plebs id quod spopondit effecit}, the common
people of the nations and of sinners did what it promised. His answer finds both sons among
John's hearers.
\latin{Primus est filius, populus ex Pharisaeis in praesens a Deo per Ioannis
prophetiam ut praeceptis suis obtemperaret admonitus}: the first son is the
people of the Pharisees, warned by God at that time through John's prophecy to
obey his precepts. Insolent and trusting in the law, it refused; but it is the
people \latin{qui deinceps operum miraculis post Domini resurrectionem poenitens
sub apostolis credidit}, which afterwards, through the miracles of the works
done after the Lord's resurrection, repented and believed under the apostles.
\latin{Filius autem minor, plebs est publicanorum et peccatorum}, the younger
son is the people of publicans and sinners, who believed John but could not yet
take up the work of the Gospel before the Passion. The Pharisees, however
unwillingly, confess that this son obeyed, \latin{obediens professione, licet
non efficiens in tempore; quia fides sola iustificat}, obedient in his
profession though not yet accomplishing it in time, because faith alone
justifies; and the chief priests and Pharisees are those \latin{qui iustificati
per fidem non erant, nec per poenitentiam regressi sunt ad salutem}, who were
not justified by faith and did not return to salvation by penance
(\work{Commentarius in Matthaeum} 21.11--15, PL 9, 1039--1041).

The disagreement is real, and it is textual as well as exegetical. For Jerome,
Chrysostom and Aquinas the son who repented is the nations; for Hilary he is
the Pharisees who believed after the resurrection. Jerome knew Hilary's answer
of v.~31 and preferred the copies that read ``the first''. Yet Hilary and
Jerome agree on what puts the publicans first, their response to John: Hilary
says they believed him, Jerome that they received his baptism of penance. And
Hilary's reading keeps in view a later faith among the very men who resisted,
the Pharisees who repented under the apostles.

\subsection{Those who are not his sheep}

The Alleluia verse is preceded in John by a harder sentence: ``you do not
believe, because you are not of my sheep'' (Jn 10:26). Augustine reads it of
those Jesus addressed without distinction: ``What did He mean, then, in saying
to them, `Ye are not of my sheep'? That He saw them predestined to everlasting
destruction, not won to eternal life by the price of His own blood.'' His
exhortation in the same passage, ``be ye sheep. They are sheep through
believing, sheep in following the Shepherd'', is addressed to his own
congregation, not to those of 10:26 (\work{Tractates on John} 48.4). On this
verse Augustine leaves the men addressed no opening.

Aquinas glosses the verse in his own terms: \latin{Vos non estis ex ovibus
meis, scilicet praedestinati ad credendum, sed praesciti ad aeternum
interitum}, you are not of my sheep, that is, not predestined to believe, but
foreknown to eternal ruin. He then asks whether saying this drove the hearers to
despair, and answers with a distinction. Common to that whole crowd was that
they \latin{non erant praeordinati a Deo ad tunc credendum}, were not
preordained by God to believe then; particular to some was that they would
believe later, \latin{aliqui qui ex eis praeordinati erant ad credendum in
posterum}, as three thousand believed in one day at Pentecost (Acts 2). Spoken
to a crowd, the saying took away no one's hope, \latin{quia nullus de se
determinate hoc poterat suspicari}, since no one could suspect it of himself in
particular; said to one determinate person, it would have (\work{Super
Ioannem} c.~10, lect.~5, p.~668). Augustine and Aquinas both read the verse of
predestination. Augustine says that those addressed were predestined to
everlasting destruction; Aquinas, that they were not predestined to believe and
were foreknown to eternal ruin, and that some of them were nevertheless
preordained to believe later. The difference is narrower than a contradiction,
and it is real: Augustine closes to those addressed the hope that Aquinas keeps
open for some of them. Aquinas's hope is for persons in that crowd, not for a
people as such.

Gregory the Great hears the chapter at another verse. Of the ``other sheep''
of 10:16 he says, \latin{Quia vero non solum Iudaeam, sed etiam gentilitatem
redimere venerat}, because the Lord had come to redeem not Judaea only but the
nations too, he spoke of bringing them; and he counts himself and his hearers
among them, \latin{qui ex gentili populo venimus}. \latin{Quasi enim ex duobus
gregibus unum ovile efficit, quia Iudaicum et gentilem populum in sua fide
coniungit}: he makes as it were one fold of two flocks, because he joins the
Jewish and the Gentile people in his faith, as Paul says, ``he is our peace,
who hath made both one'' (Eph 2:14); for \latin{ad aeternam vitam ex utraque
natione simplices eligit}, he chooses the simple from each nation for eternal
life (\work{Homiliae in Evangelia} 14.4). Gregory's fold is made of those
chosen from each people and joined in faith; he says nothing of 10:26.

\subsection{How the reading is handed on}

The identification of the second son with the Jewish people belongs to
Jerome, Chrysostom and Aquinas as their exegesis of the parable. The Second
Vatican Council's declaration \work{Nostra aetate} governs how such readings are
taught and preached today. It teaches that ``God holds the Jews most dear for
the sake of their Fathers'' and does not repent of his gifts and calls; that
what happened in the Passion ``cannot be charged against all the Jews, without
distinction, then alive, nor against the Jews of today''; and that the Jews
``should not be presented as rejected or accursed by God, as if this followed
from the Holy Scriptures'' (no.~4). The texts of the Fathers stand as they
are, and so do the limits they themselves carry. Matthew names the parable's
hearers as the chief priests and the ancients of the people (21:23) and the
chief priests and Pharisees (21:45). Jerome's sentence on Ezekiel 18:29 stands
beside his reading of 18:31 \latin{ad utrumque populum} and his new heart for
Israel. Chrysostom hears in ``go before you'' a hope, ``if they were willing''.
And Aquinas turns the parable on the Church's own members, on those who profess
and do not do.

\subsection{The Missal's texts in this reading}

In this interpretation the placing of the Missal's texts is the editor's, not a
commentator's. The Entrance
Antiphon is a confession spoken for a people. Jerome notes that the three young
men in the furnace had not sinned, and that they speak \latin{ex persona
populi}, in the person of the people (\work{In Danielem}, PL 25, 509). Heard
beside the parable, it is the late truth the second son owed: a people that
promised and did not obey confesses that God's judgments on it were true. The
Collect asks that those who run toward God's promises share the goods of
heaven, and in this reading the promises made to the fathers are shared by
those who came late.

The responsorial psalm, in Augustine's exposition, is spoken by Christ ``in the
person of the Church: for what is said has reference rather to the Christian
People turned unto God''; it is the prayer of one who, ``dismissed by Thee from
Paradise, and having taken my journey into a far country'', cannot return
``unless Thou meetest the wanderer'' (\work{Enarrationes in Psalmos} 24.1, 5).
Jerome's first remark on the parable, that its two sons are those of Luke's,
the frugal and the prodigal, sets the same far country beside this Gospel; the
join between the two expositions is the editor's.

The second reading enters this interpretation only by the editor's joining. It
is a letter to a church of the nations at Philippi, asked to be of one mind,
each esteeming the others better than himself, as the first and the second
son alike are asked to work in one vineyard. In the longer form it ends with
the day when ``every knee should bow'' and ``every tongue should confess''.
The Prayer over the Offerings asks that the offering open the fountain of every
blessing, and here the blessing promised through the fathers reaches all.

Of the two Communion antiphons, the first is heard through Hilary's
exposition of the verse. \latin{Omnis dei sermo, qui scripturis diuinis
continetur, in spem bonorum caelestium uocat}, every word of God contained in
the divine Scriptures calls to the hope of heavenly goods; and hope \latin{non
inanem esse oportet neque uerbis tantummodo praedicari, sed rebus ipsis
doceri}, must not be empty nor proclaimed in words only, but shown in the
things themselves (\work{Tractatus super Psalmos} 118, Zain~1, CSEL 22,
p.~418). Hope shown in words only is the second son's ``I go, Sir''; the join
to the parable is the editor's. The second antiphon
names the love that makes one fold. Gregory, on the shepherd who lays down his
life, writes \latin{ea charitate qua pro ovibus morior quantum Patrem diligam
ostendo}, by the charity with which I die for the sheep I show how much I love
the Father, and goes on at once to the other sheep and the one fold
(\work{Homiliae in Evangelia} 14.4); he says it of John 10, and its use for the
antiphon from John's letter is the editor's. The Prayer after Communion asks
that the communicants be fellow heirs of Christ's glory. Jerome's second son
killed the son and thought himself heir; the heirs of this prayer are those who
share the Son's sufferings. The shared vocabulary of inheritance is in the
texts, and the join is the editor's.

\subsection{Where this interpretation strains}

Hilary's contrary identification and the divided text of v.~31 show that the
allegory of the peoples is a reading of the parable, held by the three
witnesses who give it, and not its plain sense. The literal addressees are the
chief priests and ancients, and the chief priests and Pharisees, not a people
as such; Jerome's harshest sentence belongs to him, and his own commentary and
Chrysostom's set a hope beside it. Augustine, who reads Ezekiel's complaint
without naming any people, is on John 10:26 the contrary position to the hope
this reading holds out for the second son. The second reading contributes to
the interpretation only through the editor's joining.

\subsection{The four senses of this interpretation}

\begin{fourSenses}
\item[Literal.] In the temple Jesus tells the chief priests and ancients, who
would not say whether John's baptism was from heaven, that the publicans and
harlots who believed John go before them into the kingdom of God. Ezekiel
answers the house of Israel that the Lord's way is right and theirs perverse.

\item[Allegorical.] The first son is the people of the nations, who refused and
afterwards believed, and the second the people who promised at Sinai and did
not go (Jerome, Chrysostom, Aquinas); Hilary assigns them the other way, the
first to the Pharisees who believed after the resurrection and the second to
the publicans who believed John. The psalm is prayed by ``the Christian People
turned unto God'' (Augustine), and that this people is the Church of the
nations is the editor's gloss.

\item[Moral.] The baptized can become the second son: Aquinas applies the
parable to clerics and religious who profess and do not do. Those long inside
the vineyard are not to resent those who come in late, and must themselves go.

\item[Anagogical.] One fold and one shepherd, made of those chosen from each
people and joined in faith (Gregory); a new heart for Israel and for the
nations (Jerome, on Ezekiel); and, where the longer form is read, every knee
bowing and every tongue confessing that the Lord Jesus Christ is in the glory
of God the Father.
\end{fourSenses}
